\setcounter{dang}{0}
%----------------NỘI DUNG BIÊN SOẠN
% \setcounter{chapter}{5}
\chap{CÁC SỐ ĐẶC TRƯNG CỦA MẪU SỐ LIỆU KHÔNG GHÉP NHÓM}
% \setcounter{section}{11}
\section{Số gần đúng và sai số}
%-------------------------------------------
\subsection{KIẾN THỨC TRỌNG TÂM}
\subsubsection{Số gần đúng}
Trong nhiều trường hợp, ta không biết hoặc khó biết số đúng (kí hiệu là $\overline{a}$) mà chỉ tìm được giá trị khác xấp xỉ nó. Giá trị này được gọi là số gần đúng, kí hiệu là $a$.
\subsubsection{Sai số tuyệt đối và sai số tương đối}
\begin{boxdn}
	Giá trị $|a-\overline{a}|$ phản ánh mức độ sai lệch giữa số đúng $\overline{a}$ và số gần đúng $a$, được gọi là \textbf{sai số tuyệt đối} của số gần đúng $a$, kí hiệu là $\Delta_a$, tức là \fbox{$\Delta_a=|a-\overline{a}|$}
\end{boxdn}
\begin{note}
 Trên thực tế, nhiều khi ta không biết $\overline{a}$ nên cũng không biết $\Delta_a$. Tuy nhiên, ta có thể đánh giá được $\Delta_a$ không vượt quá một số dương $d$ nào đó.
\end{note}
\begin{boxdn}
Nếu biết $\Delta_a \leq d$ thì \fbox{$a-d\leq \overline{a}\leq a+d$}, khi đó ta viết \fbox{$\overline{a}=a\pm d$} và hiểu là \fbox{$\overline{a} \in [a-d;a+d]$}.\\
Do $d$ càng nhỏ thì $a$ càng gần $\overline{a}$ nên $d$ được gọi là \textbf{độ chính xác của số gần đúng}.
\end{boxdn}
\begin{boxdn}
	\textbf{Sai số tương đối} của số gần đúng $a$, kí hiệu là $\delta_a$, là tỉ số giữa sai số tuyệt đối $\Delta_a$ và $|a|$, tức là \fbox{$\delta_a=\dfrac{\Delta_a}{|a|}$}.
\end{boxdn}
\begin{nx}
	Nếu $\overline{a}=a\pm d$ thì $\Delta_a \leq d$, do đó \fbox{$\delta_a \leq \dfrac{d}{|a|}$}. Nếu $\dfrac{d}{|a|}$ càng nhỏ thì chất lượng của phép đo hay tính toán càng cao. Người ta thường viết sai số tương đối dưới dạng phần trăm.
\end{nx}
\subsubsection{Số quy tròn. Quy tròn số gần đúng}
\iconMT \indam{Số quy tròn}
\begin{boxdn}
	\begin{itemize}
		\item Trong thực tế đo đạc và tính toán, nhiều khi ta chỉ biết giá trị gần đúng của một đại lượng với độ chính xác nào đó (kể cả khi biết dược giá trị đúng của nó). Khi đó để cho gọn, các số thường được làm tròn (còn gọi là quy tròn).
		\item Số thu được sau khi thực hiện làm tròn số được gọi là số quy tròn. Số quy tròn là một số gần đúng của số ban đầu.
	\end{itemize}
\end{boxdn}
\iconMT \indam{Quy tròn số gần đúng}
\begin{nx}
	\begin{itemize}
		\item Khi thay số đúng bởi số quy tròn đến một hàng nào đó thì sai số tuyệt đối của số quy tròn không vượt quá nửa đơn vị của hàng quy tròn.
		\item Cho số gần đúng với độ chính xác $d$. Khi được yêu cầu làm tròn số $a$ mà không nói rõ đến hàng nào thì ta làm tròn số $a$ đến hàng thấp nhất mà $d$ nhỏ hơn $1$ đơn vị của hàng đó.
	\end{itemize}
\end{nx}
%-------------------------------------------
\subsection{PHÂN LOẠI VÀ PHƯƠNG PHÁP GIẢI TOÁN}
\begin{dang}{Quy tròn số gần đúng}
	\begin{itemize}
		\item Khi thay số đúng bởi số quy tròn đến một hàng nào đó thì sai số tuyệt đối của số quy tròn không vượt quá nửa đơn vị của hàng quy tròn.
		\item Cho số gần đúng với độ chính xác $d$. Khi được yêu cầu làm tròn số $a$ mà không nói rõ đến hàng nào thì ta làm tròn số $a$ đến hàng thấp nhất mà $d$ nhỏ hơn $1$ đơn vị của hàng đó.
	\end{itemize}
\end{dang}
\begin{vd}
	Cho số gần đúng $a=329\,401$ với độ chính xác $d=300$. Hãy viết số quy tròn của số $a$.
\loigiai{
	Vì độ chính xác đến hàng trăm ($d=300$) nên ta làm tròn $a$ đến hàng nghìn.\\
	Khi đó số quy tròn của $a$ là $329\,000$.
	}
\end{vd}
\begin{vd}
	Cho số gần đúng $a=329\,451$ với độ chính xác $d=20$. Hãy viết số quy tròn của số $a$.
\loigiai{
	Vì độ chính xác đến hàng chục ($d=20$) nên ta làm tròn $a$ đến hàng trăm. \\
	Khi đó số quy tròn của $a$ là $329\,500$.
	}
\end{vd}
\begin{vd}
	Cho số gần đúng $a=329{,}421$ với độ chính xác $d=0{,}01$. Hãy viết số quy tròn của số $a$.
\loigiai{
	Vì độ chính xác đến hàng phần trăm ($d=0{,}01$) nên ta làm tròn $a$ đến hàng phần chục.\\
	Khi đó số quy tròn của $a$ là $329{,}4$.
	}
\end{vd}
\begin{dang}{Vận dụng vào thực tế xác định số đúng, số gần đúng, độ chính xác của một phép đo và sai số tuyệt đối, sai số tương đối của số gần đúng}
	\begin{itemize}
		\item Khi ta viết số đúng $\overline{a}=a\pm d$, tức là $a$ là số gần đúng với số đúng $\overline{a}$ với độ chính xác $d$; đồng thời $\overline{a}$ thuộc đoạn $[a-d;a+d]$.
		\item $\Delta_a=|a-\overline{a}|$ là sai số tuyệt đối của số gần đúng $a$ và $\Delta_a\leq d$.
		\item $\delta_a=\dfrac{\Delta_a}{|a|}$ là sai số tương đối của số gần đúng $a$ và $\delta_a \leq \dfrac{d}{|a|}$.
	\end{itemize}
\end{dang}
\begin{vd}
	Một công ty sử dụng dây chuyền $A$ đề đóng gạo vào bao với khối lượng mong muốn là $20$ kg. Trên bao bì ghi thông tin khối lượng là $20\pm 0{,}2$ kg. Gọi $\overline{a}$ là khối lượng thực của một bao gạo do dây chuyền $A$ đóng gói.
	\begin{enumerate}
		\item Xác định số gần đúng và độ chính xác của phép đo khối lượng bao gạo.
		\item Giá trị thực $\overline{a}$ của bao gạo nằm trong đoạn nào?
	\end{enumerate}
\loigiai{
	\begin{enumerate}
		\item Theo đề $\overline{a}=20\pm 0{,}2$ kg nên số gần đúng cho $\overline{a}$ là $a=20$ (kg) và độ chính xác $d=0{,}2$ (kg).
		\item Giá trị thực $\overline{a}$ của bao gạo nằm trong đoạn $[20-0{,}2;20+0{,}2]$ hay $[19{,}8;20{,}2]$.
	\end{enumerate}
	}
\end{vd}
\begin{vd}
	Trong một cuộc điều tra dân số, người ta viết dân số của tỉnh $B$ là
	$$4\,271\,620 \text{ người} \pm 50\,000 \text{ người}.$$
	Hãy đánh giá sai số tương đối của số gần đúng này.
\loigiai{
	Theo đề số gần đúng $a=4\,271\,620$ người và độ chính xác $d=50\,000$ người nên sai số tương đối là
	$$\delta_a\leq \dfrac{d}{|a|}=\dfrac{50\,000}{4\,271\,620}\approx 1{,}2\%.$$
	}
\end{vd}
\begin{vd}
	Bạn Nam thực hiện thí nghiệm và xác định được khối lượng riêng của nước tinh khiết ở $4^\circ C$ là $999{,}981$ kg/m$^3$.
	\begin{enumerate}
		\item Đây là số đúng hay số gần đúng?
		\item Hãy tính sai số tuyệt đối.
	\end{enumerate}
\loigiai{
	\begin{enumerate}
		\item Giá trị $999{,}981$ kg/m$^3$ là số gần đúng cho khối lượng riêng của nước tinh khiết ở $4^\circ C$.
		\item Theo đề số đúng $\overline{a}=1\,000$ kg/m$^3$ và số gần đúng $a=999{,}981$ kg/m$^3$ nên sai số tuyệt đối là
		$$|a-\overline{a}|=1\,000-999{,}981=0{,}019.$$
	\end{enumerate}
	}
\end{vd}
%-------------------------------------------
\subsection{BÀI TẬP RÈN LUYỆN}
%-------------Trắc nghiệm nhiều lựa chọn
\ind{PHẦN I.} \inden{Câu trắc nghiệm nhiều phương án lựa chọn. Mỗi câu hỏi học sinh chỉ chọn một phương án.}
\setcounter{ex}{0}
\Opensolutionfile{ans}[ans/ans-c5b12-lc]
\begin{ex}%[0D6N1-2][Biên soạn tài liệu dạy học bộ KNTT-CS]%[Bồ Văn Hậu]
	Dùng thước đo có độ chia nhỏ nhất là $1$ cm để đo chiều cao của một nam sinh lớp $11A$. Hỏi độ chính xác của phép đo này là bao nhiêu cm?
	\choice
	{$1$ cm}
	{\True $0{,}5$ cm}
	{$4$ cm}
	{$2$ cm}
\loigiai{
	Vì độ chia nhỏ nhất của thước là $1$ cm nên độ chính xác của phép đo này là $d=\dfrac{1}{2}=0{,}5$ cm.
	}
\end{ex}
\begin{ex}%[0D6H1-3][Biên soạn tài liệu dạy học bộ KNTT-CS]%[Bồ Văn Hậu]
	Cho số gần đúng $a=689\,567$ với độ chính xác $d=100$. Hãy viết số quy tròn của số $a$.
	\choice
	{\True $690\,000$}
	{$699\,000$}
	{$689\,000$}
	{$689\,600$}
\loigiai{
	Vì độ chính xác đến hàng trăm ($d=100$) nên ta làm tròn $a$ đến hàng nghìn. \\
	Khi đó số quy tròn của $a$ là $690\,000$.
	}
\end{ex}
\begin{ex}%[0D6H1-3][Biên soạn tài liệu dạy học bộ KNTT-CS]%[Bồ Văn Hậu]
	Cho số gần đúng $a=95\,628$ với độ chính xác $d=40$. Hãy viết số quy tròn của số $a$.
	\choice
	{$95\,620$}
	{$96\,600$}
	{$96\,000$}
	{\True $95\,600$}
\loigiai{
	Vì độ chính xác đến hàng chục ($d=40$) nên ta làm tròn $a$ đến hàng trăm. \\
	Khi đó số quy tròn của $a$ là $95\,600$.
	}
\end{ex}
\begin{ex}%[0D6H1-3][Biên soạn tài liệu dạy học bộ KNTT-CS]%[Bồ Văn Hậu]
	Cho số gần đúng $a=85{,}773$ với độ chính xác $d=0{,}01$. Hãy viết số quy tròn của số $a$.
	\choice
	{$85{,}7$}
	{\True $85{,}8$}
	{$85{,}77$}
	{$85{,}87$}
\loigiai{
	Vì độ chính xác đến hàng phần trăm ($d=0{,}01$) nên ta làm tròn $a$ đến hàng phần chục. \\
	Khi đó số quy tròn của $a$ là $85{,}8$.
	}
\end{ex}
\begin{ex}%[0D6N1-2][Biên soạn tài liệu dạy học bộ KNTT-CS]%[Bồ Văn Hậu]
	Sử dụng máy tính bỏ túi hãy viết giá trị gần đúng của $\sqrt{21}$ chính xác đến hàng phần trăm.
	\choice
	{$4{,}6$}
	{\True $4{,}58$}
	{$4{,}586$}
	{$4{,}59$}
\loigiai{
	Ta có $\sqrt{21}=4{,}582575695$ nên ta làm tròn đến hàng phần trăm được kết quả là $4{,}58$.
	}
\end{ex}
\begin{ex}%[0D6N1-2][Biên soạn tài liệu dạy học bộ KNTT-CS]%[Bồ Văn Hậu]
	Sử dụng máy tính bỏ túi hãy viết giá trị gần đúng của $\sqrt[3]{2}$ chính xác đến hàng phần nghìn.
	\choice
	{$1{,}263$}
	{$1{,}2617$}
	{$1{,}259$}
	{\True $1{,}26$}
\loigiai{
	Ta có $\sqrt[3]{2}=1{,}2599210499$ nên ta làm tròn đến hàng phần nghìn được kết quả là $1{,}26$.
	}
\end{ex}
\begin{ex}%[0D6N1-2][Biên soạn tài liệu dạy học bộ KNTT-CS]%[Bồ Văn Hậu]
	Độ cao của một ngọn núi được ghi lại như sau $\overline{h}=818{,}7\pm 0{,}4$ m. Độ chính xác $d$ của phép đo trên là
	\choice
	{\True $d=0{,}4$ m}
	{$d=0{,}7$ m}
	{$d=1{,}1$ m}
	{$d=0{,}5$ m}
\loigiai{
	Theo đề $\overline{h}=818{,}7\pm 0{,}4$ m nên số gần đúng cho $\overline{h}$ là $818{,}7$ (m) và độ chính xác $d=0{,}4$ (m).
	}
\end{ex}
\begin{ex}%[0D6N1-2][Biên soạn tài liệu dạy học bộ KNTT-CS]%[Bồ Văn Hậu]
	Một công ty sử dụng dây chuyền $P$ đề đóng gạo vào bao với khối lượng mong muốn là $10$ kg. Trên bao bì ghi thông tin khối lượng là $10\pm 0{,}1$ kg. Gọi $\overline{a}$ là khối lượng thực của một bao gạo do dây chuyền $P$ đóng gói. Xác định độ chính xác của phép đo khối lượng bao gạo.
	\choice
	{$0{,}2$ kg}
	{$0{,}3$ kg}
	{$1$ kg}
	{\True $0{,}1$ kg}
\loigiai{
	Theo đề $\overline{a}=10\pm 0{,}1$ kg nên số gần đúng cho $\overline{a}$ là $10$ (kg) và độ chính xác $d=0{,}1$ (kg).
	}
\end{ex}
\begin{ex}%[0D6H1-2][Biên soạn tài liệu dạy học bộ KNTT-CS]%[Bồ Văn Hậu]
	Một công ty sử dụng dây chuyền $P$ đề đóng gạo vào bao với khối lượng mong muốn là $10$ kg. Trên bao bì ghi thông tin khối lượng là $10\pm 0{,}1$ kg. Gọi $\overline{a}$ là khối lượng thực của một bao gạo do dây chuyền $P$ đóng gói. Giá trị thực $\overline{a}$ của bao gạo nằm trong đoạn nào?
	\choice
	{\True $[9{,}9;10{,}1]$}
	{$[9{,}8;10{,}2]$}
	{$[9{,}99;10{,}01]$}
	{$[9{,}98;10{,}02]$}
\loigiai{
	Theo đề $\overline{a}=10\pm 0{,}1$ kg nên số gần đúng cho $\overline{a}$ là $a=10$ (kg) và độ chính xác $d=0{,}1$ (kg).\\
	Suy ra giá trị thực $\overline{a}$ của bao gạo nằm trong đoạn $[10-0{,}1;10+0{,}1]$ hay $[9{,}9;10{,}1]$.
	}
\end{ex}
\begin{ex}%[0D6H1-3][Biên soạn tài liệu dạy học bộ KNTT-CS]%[Bồ Văn Hậu]
	Trong một cuộc điều tra dân số, người ta viết dân số của tỉnh $C$ là
	$$7\,241\,920 \text{ người} \pm 40\,000 \text{ người}.$$
	Hỏi sai số tương đối của số gần đúng này không vượt quá số nào trong các số bên dưới?
	\choice
	{\True $0{,}55\%$}
	{$0{,}50\%$}
	{$0{,}53\%$}
	{$0{,}45\%$}
\loigiai{
	Theo đề số gần đúng $a=7\,241\,920$ người và độ chính xác $d=40\,000$ người nên sai số tương đối là
	$$\delta_a\leq \dfrac{d}{|a|}=\dfrac{40\,000}{7\,241\,920}\approx 0{,}55\%.$$
	}
\end{ex}
\begin{ex}%[0D6H1-3][Biên soạn tài liệu dạy học bộ KNTT-CS]%[Bồ Văn Hậu]
	Đo vận tốc trung bình của một chiếc xe ô tô chạy trên đường cao tốc cho ta được kết quả là
	$$98 \pm 7 \text{ km/h}.$$
	Hỏi sai số tương đối của số gần đúng này không vượt quá số nào trong các số bên dưới?
	\choice
	{$7{,}052\%$}
	{\True $7{,}143\%$}
	{$7{,}112\%$}
	{$7{,}132\%$}
\loigiai{
	Theo đề số gần đúng $a=98$ km/h và độ chính xác $d=7$ km/h nên sai số tương đối là
	$$\delta_a\leq \dfrac{d}{|a|}=\dfrac{7}{98}\approx 7{,}143\%.$$
	}
\end{ex}
\begin{ex}%[0D6H1-3][Biên soạn tài liệu dạy học bộ KNTT-CS]%[Bồ Văn Hậu]
	Bạn Minh thực hiện thí nghiệm và xác định được khối lượng riêng của nước tinh khiết ở $4^\circ C$ là $999{,}983$ kg/m$^3$. Hãy tính sai số tuyệt đối.
	\choice
	{$0{,}019$}
	{$0{,}013$}
	{\True $0{,}017$}
	{$0{,}011$}
\loigiai{
	Theo đề số đúng $\overline{a}=1\,000$ kg/m$^3$ và số gần đúng $a=999{,}983$ kg/m$^3$ nên sai số tuyệt đối là
	$$|a-\overline{a}|=1\,000-999{,}983=0{,}017.$$
	}
\end{ex}
\Closesolutionfile{ans}
\indapan{6}{ans/ans-c5b12-lc}
%---------------------Trắc nghiệm đúng sai
\ind{PHẦN II.} \inden{Câu trắc nghiệm đúng sai. Trong mỗi ý a), b), c), d) ở mỗi câu, học sinh chọn đúng hoặc sai.}
\setcounter{ex}{0}
\Opensolutionfile{ans}[ans/ans-c5b12-ds]
\begin{ex}%[0D6N1-2][Biên soạn tài liệu dạy học bộ KNTT-CS]%[Bồ Văn Hậu]
	Cho số gần đúng $a=975\,427$.
	\choiceTF
	{\True Số quy tròn của $a$ với độ chính xác $d_1=1000$ là $980\,000$}
	{Số quy tròn của $a$ với độ chính xác $d_2=200$ là $976\,000$}
	{\True Số quy tròn của $a$ với độ chính xác $d_3=40$ là $975\,400$}
	{Số quy tròn của $a$ với độ chính xác $d_4=3$ là $975\,428$}
\loigiai{
	\begin{itemchoice}
		\itemch Vì độ chính xác $d_1=1000$ nên ta làm tròn $a$ đến hàng chục nghìn, ta được kết quả là $980\,000$.
		\itemch Vì độ chính xác $d_2=200$ nên ta làm tròn $a$ đến hàng nghìn, ta được kết quả là $975\,000$.
		\itemch Vì độ chính xác $d_3=40$ nên ta làm tròn $a$ đến hàng trăm, ta được kết quả là là $975\,400$.
		\itemch Vì độ chính xác $d_4=3$ nên ta làm tròn $a$ đến hàng chục, ta được kết quả là $975\,430$.
	\end{itemchoice}
	}
\end{ex}
\begin{ex}%[0D6H1-2][Biên soạn tài liệu dạy học bộ KNTT-CS]%[Bồ Văn Hậu]
	Bác Năm đo chiều cao của một cây cà chua và ghi nhận lại như sau $\overline{a}=53\pm 2$ cm.
	\choiceTF
	{\True Số gần đúng trong phép đo trên là $a=53$ (đơn vị cm)}
	{\True Độ chính xác trong phép đo trên là $d=2$ (đơn vị cm)}
	{Chiều cao của cây cà chua trên thấp hơn $50$ cm}
	{Chiều cao thực của cây cà chua nằm trong đoạn $[55;57]$ (đơn vị cm)}
\loigiai{
	\begin{itemchoice}
		\itemch Số gần đúng trong phép đo trên là $a=53$ (đơn vị cm)
		\itemch Độ chính xác trong phép đo trên là $d=2$ (đơn vị cm)
		\itemch Vì $\overline{a}=53\pm 2$ cm nên chiều cao thực của cây cà chua nằm trên đoạn $[51;55]$ (đơn vị cm).
		\itemch Chiều cao thực của cây cà chua nằm trong đoạn $[51;55]$ (đơn vị cm)
	\end{itemchoice}
	}
\end{ex}
\begin{ex}%[0D6H1-2][Biên soạn tài liệu dạy học bộ KNTT-CS]%[Bồ Văn Hậu]
	Một công ty sử dụng dây chuyền $A$ để đóng gạo vào bao với khối lượng mong muốn $4{,}9$ kg. Trên bao bì ghi thông tin khối lượng $4{,}9 \pm 0{,}2$ kg. Gọi $\overline{a}$ là khối lượng thực tế của một bao gạo do dây chuyền $A$ đóng gói.
	\choiceTF
	{\True Độ chính xác trong phép đo trên là $0{,}2$ (đơn vị kg)}
	{\True $4{,}9$ (đơn vị kg) là số gần đúng cho $\overline{a}$ (đơn vị kg)}
	{\True Khối lượng bao gạo trên không vượt quá $5{,}1$ kg}
	{Giá trị của $\overline{a}$ nằm trong đoạn $[4{,}8; 5]$}
\loigiai{
	\begin{itemchoice}
		\itemch Độ chính xác trong phép đo trên là $d=0{,}2$ (đơn vị kg).
		\itemch Vì số đúng là $\overline{a}$ và khối lượng bao gạo là $4{,}9 \pm 0{,}2$ kg nên ta xem $4{,}9$ kg là số gần đúng cho $\overline{a}$ (đơn vị kg).
		\itemch Vì $\overline{a}=4{,}9 \pm 0{,}2$ kg nên khối lượng bao gạo nằm trên đoạn $[4{,}7; 5{,}1]$ (đơn vị kg).
		\itemch Vì khối lượng bao gạo nằm trong đoạn $[4{,}7; 5{,}1]$ (đơn vị kg).
	\end{itemchoice}
	}
\end{ex}
\begin{ex}%[0D6V1-3][Biên soạn tài liệu dạy học bộ KNTT-CS]%[Bồ Văn Hậu]
	Độ dài chiều rộng và chiều dài của mảnh vườn hình chữ nhật lần lượt được ghi nhận là $x=7{,}8$ m $\pm 2$ cm và $y=25{,}6$ m $\pm 4$ cm.
	\choiceTF
	{\True Chiều rộng của hình chữ nhật trên nằm trên đoạn $[7{,}78;7{,}82]$ (đơn vị m)}
	{\True Chiều dài của hình chữ nhật trên nằm trên đoạn $[25{,}56;25{,}64]$ (đơn vị m)}
	{\True Diện tích gần đúng của hình chữ nhật là $S=199{,}68$ m$^2$ với độ chính xác $d=0{,}8248$}
	{Diện tích (sau khi quy tròn) của mảnh vườn hình chữ nhật trên lớn hơn $200$ m$^2$}
\loigiai{
	\begin{itemchoice}
		\itemch Chiều rộng của hình chữ nhật là $x=7{,}8$ m $\pm 2$ cm nên sẽ nằm trên đoạn $[7{,}78;7{,}82]$ (đơn vị m).
		\itemch Chiều dài của hình chữ nhật là $y=25{,}6$ m $\pm 4$ cm nên sẽ nằm trên đoạn $[25{,}56;25{,}64]$ (đơn vị m).
		\itemch Ta có $\heva{&7{,}78\leq x\leq 7{,}82\\&25{,}56\leq y\leq 25{,}64}\Rightarrow 198{,}8568\leq \overline{S}=xy\leq 200{,}5048$.\\
		Diện tích gần đúng là $S=7{,}8\cdot 25{,}6=199{,}68$.\\
		Suy ra $-0{,}8232\leq \overline{S} - S\leq 0{,}8248\Rightarrow\Delta_S=|\overline{S} - S|\leq 0{,}8248$.\\
		Suy ra diện tích gần đúng là $S=199{,}68$ m$^2$ với độ chính xác $d=0{,}8248$.
		\itemch Số quy tròn của $S$ là $200$ m$^2$.
	\end{itemchoice}
	}
\end{ex}
\Closesolutionfile{ans}
\indapan{3}{ans/ans-c5b12-ds}
%---------------Trắc nghiệm trả lời ngắn
\ind{PHẦN III.} \inden{Câu trả lời ngắn.}
\setcounter{ex}{0}
\Opensolutionfile{ans}[ans/ans-c5b12-kq]
%-----Câu 1
\begin{ex}%[0D6H1-3]%[Biên soạn tài liệu dạy học bộ KNTT-CS]%[Bồ Văn Hậu]
	Cho số gần đúng $a=2\,927$ với độ chính xác $d=17$. Hãy viết số quy tròn của số $a$.
	\par
	\shortans{2900}
\loigiai{
	Vì độ chính xác đến hàng chục ($d=17$) nên ta làm tròn $a$ đến hàng trăm. \\
	Khi đó số quy tròn của $a$ là $2\,900$.
	}
\end{ex}
\begin{ex}%[0D6H1-3]%[Biên soạn tài liệu dạy học bộ KNTT-CS]%[Bồ Văn Hậu]
	Cho số gần đúng $a=1{,}926$ với độ chính xác $d=0{,}01$. Hãy viết số quy tròn của số $a$.
	\par
	\shortans{1,9}
\loigiai{
	Vì độ chính xác đến hàng phần trăm ($d=0{,}01$) nên ta làm tròn $a$ đến hàng phần chục.\\
	Khi đó số quy tròn của $a$ là $1{,}9$.
	}
\end{ex}
\begin{ex}%[0D6H1-3]%[Biên soạn tài liệu dạy học bộ KNTT-CS]%[Bồ Văn Hậu]
	Tìm số gần đúng đến hàng phần trăm của số $3+2\sqrt{5}$.
	\par
	\shortans{7,47}
\loigiai{
	Ta có $3+2\sqrt{5}=7{,}472135955\ldots$ nên số gần đúng đến hàng phần trăm của số $3+2\sqrt{5}$ là $7{,}47$.
	}
\end{ex}
\begin{ex}%[0D6H1-2]%[Biên soạn tài liệu dạy học bộ KNTT-CS]%[Bồ Văn Hậu]
	Độ cao của núi Bà Đen tỉnh Tây Ninh được ghi lại như sau $\overline{h}=986\pm 0{,}3$ m. Gọi độ chính xác của phép đo trên là $d$. Giá trị của $70d-62$ là bao nhiêu?
	\par
	\shortans{-41}
\loigiai{
	Theo đề $\overline{h}=986\pm 0{,}3$ m nên số gần đúng cho $\overline{h}$ là $986$ (m) và độ chính xác $d=0{,}3$ (m).\\
	Vậy $70d-62=70\cdot 0{,}3-62=-41$.
	}
\end{ex}
\begin{ex}%[0D6H1-2]%[Biên soạn tài liệu dạy học bộ KNTT-CS]%[Bồ Văn Hậu]
	Một công ty sử dụng dây chuyền $K$ đề đóng gói chè vào bao bì với khối lượng mong muốn là $200$ gam. Trên bao bì ghi thông tin khối lượng là $200 \pm 3$ gam. Gọi $\overline{a}$ là khối lượng thực của một bao bì chè do dây chuyền $K$ đóng gói. Giá trị thực $\overline{a}$ của bao bì chè nằm trong đoạn $[m;n]$. Giá trị của biểu thức $2m-n$ là bao nhiêu?
	\par
	\shortans{191}
\loigiai{
	\begin{itemize}
		\item Theo đề $\overline{a}=200 \pm 3$ gam nên số gần đúng cho $\overline{a}$ là $a=200$ (gam) và độ chính xác $d=3$ (gam).
		\item Giá trị thực $\overline{a}$ của bao bì chè nằm trong đoạn $[200-3;200+3]$ hay $[197;203]$.
	\end{itemize}
	Vậy $m=197$, $n=203$, suy ra $2m-n=2\cdot 197-203=191$.
	}
\end{ex}
\begin{ex}%[0D6V1-1]%[Biên soạn tài liệu dạy học bộ KNTT-CS]%[Bồ Văn Hậu]
	Bác nông dân đo mảnh vườn hình chữ nhật có chiều dài $ 4{,}79\pm 0{,}05$ m và chiều rộng $ 4{,}39\pm 0{,}02 $ m. Sai số tương đối của phép đo diện tích mảnh vườn không vượt quá $a\%$ (kết quả làm tròn đến hàng phần chục). Giá trị của $a^2+5$ là bao nhiêu?
	\par
	\shortans[0]{7{,}25}
\loigiai{
	Gọi $x$, $y$ lần lượt là chiều dài và chiều rộng thực của mảnh vườn (đơn vị m).\\
	Ta có $$\heva{&x=4{,}79\pm0.05\text{ m}\\ &y=4{,}39\pm0.02\text{ m}} \Rightarrow \heva{&4{,}74\leq x\leq4{,}84\\ &4{,}37\leq y\leq4{,}41.}$$
	Gọi diện tích mảnh vườn là S, khi đó $ 20{,}7487\leq S\leq21{,}3798 \Rightarrow S=21{,}06\pm0{,}32 $ m$ ^2 $.\\
	Suy ra sai số tương đối trong phép đo là $ \delta\leq\dfrac{0{,}32}{21{,}06} \approx 0{,}015= 1{,}5\% $.\\
	Vậy $a=1{,}5$, suy ra $a^2+5=7{,}25$.
	}
\end{ex}
\Closesolutionfile{ans}
\indapan{4}{ans/ans-c5b12-kq}
%------------------Tự luận
\ind{PHẦN IV.} \inden{Tự luận.}
\setcounter{ex}{0}
\begin{ex}%[0D6H1-3]%[Biên soạn tài liệu dạy học bộ KNTT-CS]%[Bồ Văn Hậu]
	Cho số gần đúng $a=145\,927$ với độ chính xác $d=37$. Hãy viết số quy tròn của số $a$.
\loigiai{
	Vì độ chính xác đến hàng chục ($d=37$) nên ta làm tròn $a$ đến hàng trăm. Khi đó số quy tròn của $a$ là $145\,900$.
	}
\end{ex}
\begin{ex}%[0D6H1-3]%[Biên soạn tài liệu dạy học bộ KNTT-CS]%[Bồ Văn Hậu]
	Cho số gần đúng $a=123{,}4567$ với độ chính xác $d=0{,}001$. Hãy viết số quy tròn của số $a$.
\loigiai{
	Vì độ chính xác đến hàng phần nghìn ($d=0{,}001$) nên ta làm tròn $a$ đến hàng phần trăm. Khi đó số quy tròn của $a$ là $123{,}46$.
	}
\end{ex}
\begin{ex}%[0D6H1-1]%[Biên soạn tài liệu dạy học bộ KNTT-CS]%[Bồ Văn Hậu]
	Trong một cuộc điều tra về số khách ghé tham quan khu du lịch Núi Bà Đen tỉnh Tây Ninh trong tháng 01/2026 của Chi cục Thống kê tỉnh Tây Ninh ghi nhận được có số khách ghé tham quan là
	$$4\,628\,794 \text{ người} \pm 3\,000 \text{ người}.$$
	Hãy đánh giá sai số tương đối của số gần đúng này.
\loigiai{
	Theo đề số gần đúng $a=4\,628\,794$ người và độ chính xác $d=3\,000$ người nên sai số tương đối là
	$$\delta_a\leq \dfrac{d}{|a|}=\dfrac{3\,000}{4\,628\,794}\approx 0{,}065\%.$$
	}
\end{ex}
\begin{ex}%[0D6H1-2]%[Biên soạn tài liệu dạy học bộ KNTT-CS]%[Bồ Văn Hậu]
	Một công ty sử dụng dây chuyền $F$ đề đóng gói cà phê vào bao bì với khối lượng mong muốn là $2$ kg. Trên bao bì ghi thông tin khối lượng là $2\pm 0{,}02$ kg. Gọi $\overline{a}$ là khối lượng thực của một bao bì cà phê do dây chuyền $F$ đóng gói.
	\begin{enumerate}
		\item Xác định số gần đúng và độ chính xác của phép đo khối lượng bao bì cà phê.
		\item Giá trị thực $\overline{a}$ của bao bì cà phê nằm trong đoạn nào?
	\end{enumerate}
\loigiai{
	\begin{enumerate}
		\item Theo đề $\overline{a}=2\pm 0{,}02$ kg nên số gần đúng cho $\overline{a}$ là $2$ (kg) và độ chính xác $d=0{,}02$ (kg).
		\item Giá trị thực $\overline{a}$ của bao bì cà phê nằm trong đoạn $[2-0{,}02;2+0{,}02]$ hay $[1{,}98;2{,}02]$.
	\end{enumerate}
	}
\end{ex}
\begin{ex}%[0D6H1-1]%[Biên soạn tài liệu dạy học bộ KNTT-CS]%[Bồ Văn Hậu]
	Bạn Linh thực hiện thí nghiệm và xác định được khối lượng riêng của nước tinh khiết ở $10^\circ C$ là $998{,}785$ kg/m$^3$.
	\begin{enumerate}
		\item Đây là số đúng hay số gần đúng?
		\item Hãy tính sai số tuyệt đối.
	\end{enumerate}
\loigiai{
	\begin{enumerate}
		\item Giá trị $998{,}785$ kg/m$^3$ là số gần đúng cho khối lượng riêng của nước tinh khiết ở $10^\circ C$.
		\item Theo đề số đúng $\overline{a}=1\,000$ kg/m$^3$ và số gần đúng $a=998{,}785$ kg/m$^3$ nên sai số tuyệt đối là
		$$|a-\overline{a}|=1\,000-998{,}785=1{,}215.$$
	\end{enumerate}
	}
\end{ex}
\begin{ex}%[0D6V1-2]%[Biên soạn tài liệu dạy học bộ KNTT-CS]%[Bồ Văn Hậu]
	Bạn An tính diện tích của hình tròn bán kính $r=4$ cm bằng công thức $S=3{,}145 \cdot 4^2=50{,}32$ cm$^2$. Biết rằng $3{,}14<\pi<3{,}15$, hãy ước lượng độ chính xác của $S$.
\loigiai{
	Gọi $\overline{S}$ là diện tích đúng của hình tròn trên.\\
	Khi đó
	\[3{,}14\cdot4^2<\overline{S}<3{,}15\cdot4^2 \Leftrightarrow 50{,}24<\overline{S}<50{,}4.
	\]
	Do đó $50{,}24-50{,}32=-0{,}08<\overline{S}-S<50{,}4-50{,}32=0{,}08$, tức là $\left| S-\overline{S}\right| <0{,}08$.\\
	Vậy kết quả của bạn An có độ chính xác là $0{,}08$.
	}
\end{ex}
\Closesolutionfile{ans}